\uuid{H8AW}
\titre{Intégration par parties}
\auteur{hanani ; mzari}
\organisation{exo7}
\licence{CC BY-NC-SA 4.0}
\datecreate{2020-01-20}
\niveau{L1}
\module{Analyse}
\chapitre{Calcul d'intégrales}
\sousChapitre{Primitives diverses}
\difficulte{3}
\theme{Primitives, intégration par parties}
\exo7qcmid{lille-2.primitives.034}

\contenu{
\begin{question}
\qclassification{127.11}
\qtags{motcle=intégration par parties}
On note par $F$ une primitive de $f(x)=\arctan (x)$ sur $\Rr$. Parmi les propositions suivantes, cocher celles qui sont vraies :
\begin{answers}
\bad{$\displaystyle F(x)=\frac{1}{1+x^2}+k$, où $k\in \Rr$.}
\bad{$\displaystyle F(x)=x\arctan (x)+k$, où $k\in \Rr$.}
\good{$\displaystyle F(x)=x\arctan (x)-\int \frac{x\mathrm{d}x}{1+x^2}$.}
\good{$\displaystyle F(x)=x\arctan (x)-\frac{1}{2}\ln (1+x^2)+k$, où $k\in \Rr$.}
\end{answers}
\begin{explanations}
Une intégration par parties avec $u=\arctan (x)$ et $v=x$ donne
$$\int \arctan (x)\mathrm{d}x=x\arctan (x)-\int \frac{x\mathrm{d}x}{1+x^2}=x\arctan (x)-\frac{1}{2}\ln (1+x^2)+k,\; k\in \Rr.$$
\end{explanations}
\end{question}
}
